% !TEX root = ../main.tex
\chapter{AURORA: ferramentas de produtividade}
\label{cap:13-aurora-produtividade}

A \tool{AURORA} reúne, em torno do editor \tool{Monaco}, um conjunto de
recursos auxiliares que aproximam a IDE da experiência de um ambiente
profissional de desenvolvimento: um terminal de saída rico, controle de versão
embutido, busca global no projeto, servidores de linguagem (LSP) para
\textit{Verilog}/\textit{SystemVerilog}, realce sintático por
\textit{tree-sitter}, formatação de código, paleta de comandos,
internacionalização, painel de configurações, atualizador automático e um
\textit{overlay} de diagnóstico de desempenho para desenvolvimento.

Este capítulo descreve cada um desses recursos a partir do código-fonte: o que
faz, qual ferramenta o sustenta e como está ligado entre o processo principal
(\textit{main}) do \tool{Electron} e o processo de renderização (\textit{renderer}).
O fluxo de compilação propriamente dito é tratado em outros capítulos; aqui o
foco é na infraestrutura de produtividade que cerca a edição.

\begin{nota}
Um padrão recorrente atravessa todos os recursos: \textbf{best-effort}. Cada
módulo verifica se o binário ou recurso de que depende está presente e, quando
não está, vira um \textit{no-op} silencioso — o editor continua funcionando
exatamente como antes, sem erro visível ao usuário. Isso aparece em
\file{verible\_lsp.js}, \file{slang\_lsp.js}, \file{clang\_format.js},
\file{grammars.js} e no atualizador.
\end{nota}

\section{Arquitetura de ligação: \textit{main} $\leftrightarrow$ \textit{renderer}}
\label{sec:13-arquitetura-ligacao}

Quase todos os recursos deste capítulo seguem a mesma topologia em três
camadas:

\begin{enumerate}[leftmargin=*]
  \item Um módulo no processo \textit{main} (sob \path{aurora/main/}) que faz o
        trabalho privilegiado: gerar processos filhos, ler o sistema de
        arquivos, falar com a API do GitHub. Ele registra \textit{handlers} de
        IPC via \lstinline|ipcMain.handle(...)|.
  \item Uma ponte exposta no \file{preload.js} (em \path{js/app/preload.js})
        via \lstinline|contextBridge.exposeInMainWorld(...)|, que publica um
        objeto no \lstinline|window| do \textit{renderer} (por exemplo
        \lstinline|window.lspAPI|, \lstinline|window.gitAPI|,
        \lstinline|window.treeSitterAPI|, \lstinline|window.clangFormatAPI| e
        \lstinline|window.slangAPI|). Algumas operações não têm ponte própria e
        são agregadas em \lstinline|window.electronAPI| (por exemplo a busca, via
        \lstinline|electronAPI.searchInProject|).
  \item Um módulo no \textit{renderer} (sob \path{aurora/js/}) que consome essa
        ponte e a integra ao \tool{Monaco} ou à interface.
\end{enumerate}

A \autoref{tab:13-mapa-modulos} resume o mapeamento entre recurso, módulo
\textit{main}, ponte e módulo \textit{renderer}.

\begin{table}[h]
\centering
\caption{Mapa de módulos das ferramentas de produtividade.}
\label{tab:13-mapa-modulos}
\small
\begin{tabularx}{\linewidth}{Y Y Y}
\toprule
\textbf{Recurso} & \textbf{Módulo \textit{main}} & \textbf{Renderer / ponte} \\
\midrule
Terminal de saída    & ---                                & \file{js/terminal/terminal\_module.js} \\
Source control       & \file{main/ipc/git.js}             & \file{js/git/git\_panel.js} (\lstinline|gitAPI|) \\
Auth GitHub          & \file{main/ipc/github\_auth.js}    & \file{js/git/git\_panel.js} \\
Busca no projeto     & \file{main/ipc/search.js}          & \file{js/search/search\_panel.js} \\
LSP Verilog          & \file{main/lsp/verible\_lsp.js}    & \file{js/editor/lsp\_integration.js} (\lstinline|lspAPI|) \\
LSP semântico        & \file{main/lsp/slang\_lsp.js}      & \file{js/editor/slang\_integration.js} (\lstinline|slangAPI|) \\
Realce tree-sitter   & \file{main/treesitter/grammars.js} & \file{js/editor/treesitter\_highlight.js} (\lstinline|treeSitterAPI|) \\
Formatação C/C++     & \file{main/format/clang\_format.js}& \file{js/editor/clang\_format\_integration.js} \\
Paleta de comandos   & ---                                & \file{js/ui/command\_palette.js} \\
i18n                 & ---                                & \file{js/i18n/i18n.js} + \path{locales/} \\
Configurações        & ---                                & \file{js/processors/aurora\_settings.js} \\
Atualizador          & \file{main/updater.js}             & \path{html/update-notification.html} \\
Overlay de jank      & ---                                & \file{js/dev/jank\_overlay.js} \\
\bottomrule
\end{tabularx}
\end{table}

\section{Terminal de saída}
\label{sec:13-terminal}

O painel inferior da \tool{AURORA} é um console de \textit{scrollback} rico,
implementado pela classe \lstinline|TerminalManager| em
\file{js/terminal/terminal\_module.js}. Não é um terminal interativo (não há
\textit{shell}); é uma superfície de \textbf{saída} estruturada que recebe e
classifica as linhas emitidas pelas ferramentas do \textit{toolchain}.

\subsection{Abas e troca de terminal}

O painel é dividido em abas, cada uma com seu próprio corpo de rolagem,
mapeadas no construtor de \lstinline|TerminalManager|:

\begin{table}[h]
\centering
\caption{Abas do terminal de saída.}
\label{tab:13-terminal-abas}
\small
\begin{tabularx}{\linewidth}{Y Y}
\toprule
\textbf{Identificador} & \textbf{Conteúdo} \\
\midrule
\lstinline|tcmm|   & Compilação \cmm \\
\lstinline|tasm|   & Montagem (\textit{assembly}) \\
\lstinline|tveri|  & Síntese \textit{Verilog} / \tool{PRISM} \\
\lstinline|twave|  & Geração de \textit{waveform} \\
\lstinline|thtest| & \textit{Terminal Hardware Test} — etapas e barra de progresso \\
\lstinline|tcmd|   & Saída de comandos genéricos \\
\bottomrule
\end{tabularx}
\end{table}

A troca de aba é feita por \file{js/terminal/terminal.js}, que exporta a função
\lstinline|switchTerminal(targetId)|. Essa é a única implementação da troca
(consolidada de cópias que existiam em \file{compilation\_flow.js} e
\file{wave\_config\_manager.js}). Ela esconde todos os corpos
(\lstinline|.terminal-content|), revela o alvo e desliza um único indicador de
aba ativa (\lstinline|.terminal-tab-indicator|) via \lstinline|transform:
translateX(...)|, em vez de ligar/desligar uma barra por aba. Os botões de
compilação focam automaticamente a aba onde sua saída cai (por exemplo,
\lstinline|cmmcomp| foca \lstinline|terminal-tcmm|).

\subsection{Cartões por tipo de mensagem}

A saída não é texto cru: cada mensagem é classificada e agrupada em
\textbf{cartões} (\lstinline|.log-entry|). O \lstinline|TerminalManager| mantém
contadores por terminal nas categorias \lstinline|error|, \lstinline|warning|,
\lstinline|success| e \lstinline|tips|, exibidos como \textit{badges} sobre os
botões de filtro (\lstinline|createCounterBadges| /
\lstinline|updateCounterDisplay|).

Mensagens do mesmo tipo emitidas na mesma sessão são consolidadas num único
cartão agrupado (\lstinline|createGroupedCard| / \lstinline|addMessageToCard|):
o cartão recebe um \textit{timestamp} em formato \lstinline|pt-BR| e cada nova
linha vira um \lstinline|.grouped-message| dentro do contêiner. Botões de
filtro permitem mostrar/ocultar categorias (\lstinline|_cardMatchesFilter|,
\lstinline|applyFilter|); o filtro trata \lstinline|tips| e \lstinline|info|
como equivalentes. Há ainda um modo \textit{verbose} (persistido em
\lstinline|localStorage| sob \lstinline|terminal-verbose-mode|) que controla a
exibição das linhas \lstinline|plain| (não classificadas).

\begin{nota}
Para manter o consumo de memória e o desempenho de rolagem limitados num
\textit{build} que despeja milhares de linhas, há dois tetos: no máximo
\lstinline|MAX_TERMINAL_ENTRIES| (5000) nós \lstinline|.log-entry| por terminal
e no máximo \lstinline|MAX_GROUPED_MESSAGES| (5000) sub-mensagens por cartão
agrupado. Os mais antigos são descartados do topo, e os contadores são
recalculados a partir do DOM já podado.
\end{nota}

\subsection{Números de linha clicáveis}

O método \lstinline|makeLineNumbersClickable(text)| transforma referências de
diagnóstico em \textit{links} navegáveis. Há dois formatos reconhecidos:

\begin{description}[leftmargin=*]
  \item[Estilo \textit{toolchain} C] padrão \lstinline|<arquivo>:<linha>:|
        emitido por \tool{iverilog}, \tool{yosys}, \tool{gcc} e pelo
        \textit{appcomp}. A expressão regular captura o caminho (inclusive um
        prefixo de unidade do Windows como \lstinline|C:\foo\bar.v:15:|) e a
        linha; o caminho é guardado em \lstinline|data-file| para que o clique
        abra exatamente aquele arquivo.
  \item[Estilo \tool{AURORA}/\tool{YANC}] padrão \enquote{\lstinline|linha N|}
        ou \enquote{\lstinline|line N|}, sem prefixo de arquivo. Nesse caso o
        clique recai sobre o último \cmm{} compilado, resolvido pelo gerenciador
        de compilação.
\end{description}

Cada \textit{link} (\lstinline|.line-link|) tem seu clique conectado por
\lstinline|_attachLineLinkClicks|, e a navegação até a posição é feita por
\lstinline|goToLine(lineNumber)|, que comanda o \tool{Monaco} a posicionar e
revelar a linha-alvo (limitada ao número total de linhas do modelo). Há também
um \textit{link} de pasta (\lstinline|appendFolderLink|), construído com
\lstinline|textContent| (sem \lstinline|innerHTML|), que alterna a árvore de
arquivos para a visão de pastas e revela a pasta correspondente via
\lstinline|standardTreeRenderer.revealFolder|.

\subsection{Barra de progresso de teste de \textit{hardware}}

O fluxo de \textit{Terminal Hardware Test} (aba \lstinline|thtest|) usa
\lstinline|renderHardwareProgress(terminalId, p)| para exibir uma barra de
progresso real em DOM. Diferentemente dos cartões, esse é um \textbf{único}
elemento \lstinline|.hw-progress| que se atualiza no lugar: ele \textbf{não} é
um \lstinline|.log-entry|, portanto o filtro de \textit{verbose} e os
contadores o ignoram e ele permanece sempre visível. A cada atualização o nó é
reanexado ao fim do terminal, para acompanhar a última saída transmitida. O
preenchimento (\lstinline|.hw-progress-fill|) anima por
\lstinline|transform: scaleX(...)|, e a linha de metadados mostra
\lstinline|ciclo/total| mais o total de leituras de entrada (rotulado via
i18n por \lstinline|terminal.htest.reads|). O \textit{payload} \lstinline|p|
carrega \lstinline|pct|, \lstinline|cyc|, \lstinline|total|, \lstinline|reads|,
\lstinline|label| e \lstinline|done|.

\section{Controle de versão (Source Control) embutido}
\label{sec:13-git}

A \tool{AURORA} embute um painel de controle de versão completo. O lado
\textit{main} vive em \file{main/ipc/git.js} e é construído sobre a biblioteca
\textbf{simple-git}, um invólucro fino sobre o binário nativo \lstinline|git|.
Como ele dirige o \lstinline|git| real, \path{.gitignore}, \textit{diffs},
\textit{merges} e credenciais se comportam exatamente como na linha de comando.
A interface vive em \file{js/git/git\_panel.js}.

\subsection{Escopo das operações}

Todas as operações agem sobre o diretório do \textbf{projeto aberto}, derivado
de \lstinline|state.currentOpenProjectPath| (a fonte única de verdade): a função
\lstinline|projectDir()| devolve o diretório-pai do arquivo \path{.spf}. As
\textit{handlers} de leitura podem honrar um diretório explícito
(\lstinline|opts.dir|) para navegar um repositório recém-clonado que ainda não
tem \path{.spf} aberto; as mutações sempre agem sobre o projeto aberto. Toda
\textit{handler} é embrulhada por \lstinline|safe(fn)|, que converte exceções em
um objeto \lstinline|{ ok:false, error }| em vez de propagá-las pelo IPC.

\subsection{Inspeção e \textit{diffs}}

As \textit{handlers} de leitura cobrem o conjunto típico de inspeção:

\begin{longtable}{>{\raggedright\arraybackslash}p{0.34\linewidth} >{\raggedright\arraybackslash}p{0.6\linewidth}}
\caption{Handlers de inspeção do controle de versão.}\label{tab:13-git-read}\\
\toprule
\textbf{Canal IPC} & \textbf{Função} \\
\midrule
\endfirsthead
\toprule
\textbf{Canal IPC} & \textbf{Função} \\
\midrule
\endhead
\lstinline|git:is-repo|      & Verifica se o diretório é um repositório. \\
\lstinline|git:status|       & \textit{Status} com \textit{branch}, \textit{tracking}, \textit{ahead}/\textit{behind} e \texttt{flags} por arquivo; opcionalmente anexa contagens \texttt{+}/\texttt{-}. \\
\lstinline|git:diff|         & \textit{Diff} unificado de um arquivo ou de toda a árvore de trabalho. \\
\lstinline|git:commit-files| & Lista (apenas \textit{numstat}) os arquivos tocados por um \textit{commit}. \\
\lstinline|git:show|         & \textit{Diff} de um arquivo dentro de um \textit{commit} (carregado sob demanda). \\
\lstinline|git:log|          & Histórico de \textit{commits} (padrão: 50). \\
\lstinline|git:branches|     & \textit{Branches} locais e remotas (\lstinline|-a|), com remotas-só destacadas. \\
\lstinline|git:remotes|      & Lista de \textit{remotes}. \\
\lstinline|git:info|         & Nome derivado de \lstinline|origin| (\lstinline|owner/repo|) e URL. \\
\bottomrule
\end{longtable}

O texto dos \textit{diffs} é limitado por \lstinline|capDiff|: um \textit{diff}
de \textit{commit} pode ter megabytes (dados \path{.mif}/\path{.hex} gerados,
\textit{blobs} versionados), então a saída é cortada num limite de linha em
\lstinline|MAX_DIFF_BYTES| ($\approx$ 600 KB) e marcada como truncada. As
contagens por arquivo da árvore de trabalho (\lstinline|workNumstat|) combinam
dois \textit{diffs} \lstinline|--numstat| (\textit{staged} + \textit{unstaged}).
Na interface, os \textit{diffs} são renderizados por \textbf{diff2html}
(importado em \file{js/git/git\_panel.js} como \lstinline|html as renderDiff|),
com a CSS \path{diff2html/bundles/css/diff2html.min.css}, em modo
\lstinline|line-by-line| e \lstinline|colorScheme: 'dark'|. Há um teto adicional
de linhas no lado do \textit{renderer} antes de entregar o conteúdo ao
\lstinline|diff2html|.

\subsection{Mutações e operações remotas}

As mutações cobrem o ciclo completo de trabalho: \lstinline|git:init|,
\lstinline|git:stage| / \lstinline|git:stage-all| / \lstinline|git:unstage|,
\lstinline|git:discard|, \lstinline|git:commit| (com \lstinline|--amend|
opcional), \lstinline|git:undo-last-commit| (\textit{reset} \textit{soft}),
\lstinline|git:checkout| (incluindo criar \textit{branch} local e rastrear uma
remota com \lstinline|--track|), \lstinline|git:merge| (\lstinline|--no-edit|),
e o conjunto de \textit{stash} (\lstinline|git:stash|, \lstinline|git:stash-list|,
\lstinline|git:stash-pop|, \lstinline|git:stash-drop|), de modo que trocar de
\textit{branch} com a árvore suja simplesmente funciona.

As operações remotas — \lstinline|git:fetch|, \lstinline|git:pull|,
\lstinline|git:push| — usam o invólucro \lstinline|remoteGit()|. O
\lstinline|pull| roda com \lstinline|--no-edit --autostash| (nenhum
\lstinline|$EDITOR| abre e mudanças locais não impedem o \textit{merge}); o
\lstinline|push| define \textit{upstream} (\lstinline|-u origin <branch>|)
apenas quando ainda não há rastreamento. O \lstinline|git:clone| transmite
progresso ao vivo: o \textit{plugin} de progresso do \textbf{simple-git} analisa
o \texttt{stderr} do \lstinline|git --progress| e cada \textit{tick} é
encaminhado ao \textit{renderer} pelo canal \lstinline|git:clone-progress|, que
move a barra de progresso do painel. Após um \textit{clone}, \lstinline|git:scan-spf|
varre o diretório (profundidade $\le$ 5) em busca de arquivos \path{.spf} para
abrir o projeto recém-clonado.

\subsection{Autenticação GitHub}
\label{sec:13-github-auth}

O \enquote{conectar sua conta GitHub} vive em \file{main/ipc/github\_auth.js} e
oferece dois caminhos de autenticação:

\begin{enumerate}[leftmargin=*]
  \item \textbf{\textit{Personal Access Token} (PAT)} — o usuário cola um
        \textit{token} (clássico ou \textit{fine-grained} com escopo
        \lstinline|repo|); ele é validado contra \lstinline|/user| da API do
        GitHub e então armazenado.
  \item \textbf{OAuth \textit{Device Flow}} — \lstinline|deviceFlowLogin| pede
        um \textit{device}/\textit{user code}, abre a página de verificação no
        navegador (\lstinline|shell.openExternal|) e faz \textit{polling} até a
        autorização. O \textit{Client ID} é público (o \textit{device flow} não
        usa segredo) e os escopos são \lstinline|repo read:org|. Ao autorizar, a
        janela da \tool{AURORA} é trazida de volta ao primeiro plano.
\end{enumerate}

Em ambos os casos o \textit{token} é cifrado com
\lstinline|safeStorage.encryptString| (DPAPI no Windows) e gravado em
\path{userData/aurora-github.json}; o texto plano nunca toca o disco. O
\textit{token} é lido somente no lado \textit{main} — o \file{git.js} o injeta
como cabeçalho \lstinline|http.extraHeader| de uso único (\lstinline|-c|, nunca
escrito na configuração do repositório) para \textit{push}/\textit{pull}/
\textit{clone}; o \textit{renderer} só consegue perguntar \enquote{quem está
conectado?} (\lstinline|github:status|), nunca obter os \textit{bytes}. Outras
\textit{handlers}: \lstinline|github:create-repo| (cria repositório sob a conta),
\lstinline|github:list-repos| (lista repositórios próprios, de colaboração e de
organização, paginando) e \lstinline|github:disconnect|. O \textit{avatar} é
baixado e convertido em \textit{data URL} para passar pela CSP do
\textit{renderer} sem afrouxá-la.

\section{Busca no projeto (Find in Files)}
\label{sec:13-busca}

A busca global, equivalente ao painel \textit{Search} do VS Code, é
implementada em \file{main/ipc/search.js} e exposta pelo canal
\lstinline|search:in-project|. O \textit{renderer} é \file{js/search/search\_panel.js},
que agrupa os resultados por arquivo (cabeçalho recolhível + linhas de
ocorrência) e abre o arquivo na linha clicada via \lstinline|TabManager|.

\begin{nota}
O motor de busca \textbf{não é o ripgrep nem um binário externo}. É uma
varredura recursiva síncrona e \textbf{sem dependências}, escrita diretamente
sobre o módulo \lstinline|fs| do Node, enraizada no diretório do projeto aberto.
\end{nota}

\subsection{Consulta e modificadores}

O \textit{payload} aceita \lstinline|{ query, caseSensitive?, wholeWord?, regex? }|.
A função \lstinline|buildRegex| monta a expressão regular: em modo literal o
texto é escapado por \lstinline|escapeRegExp|; \lstinline|wholeWord| envolve o
corpo em \lstinline|\b...\b|; \lstinline|caseSensitive| controla a \textit{flag}
\lstinline|i|. Um padrão inválido em modo \textit{regex} é convertido em
\lstinline|{ ok:false, error }|. Os modificadores de caixa/palavra/\textit{regex}
são persistidos no \textit{renderer} sob \lstinline|aurora.search.toggles|.

\subsection{Limites da varredura}

Como a varredura é síncrona, há limites rígidos que impedem que uma árvore
grande congele o processo principal:

\begin{table}[h]
\centering
\caption{Limites da busca em \file{main/ipc/search.js}.}
\label{tab:13-busca-limites}
\small
\begin{tabularx}{\linewidth}{Y Y}
\toprule
\textbf{Constante} & \textbf{Valor / efeito} \\
\midrule
\lstinline|MAX_DEPTH|         & 12 níveis de profundidade. \\
\lstinline|MAX_FILE_BYTES|    & $\approx$ 1,5 MB — arquivos maiores são ignorados. \\
\lstinline|BINARY_SNIFF_BYTES|& 4 KB lidos para detectar byte \texttt{NUL} (binário). \\
\lstinline|PREVIEW_MAX|       & 240 caracteres por linha de \textit{preview}. \\
\lstinline|MAX_MATCHES|       & 2000 ocorrências no total. \\
\lstinline|MAX_FILES|         & 500 arquivos com ocorrência. \\
\bottomrule
\end{tabularx}
\end{table}

Diretórios de VCS, dependências e \textit{build} nunca são percorridos
(\lstinline|SKIP_DIRS|: \path{.git}, \path{node_modules}, \path{dist},
\path{build}, \path{components}, \path{.vite} entre outros). Arquivos binários
são detectados por uma heurística simples (\lstinline|looksBinary|: byte
\texttt{NUL} nos primeiros 4 KB). Ao atingir qualquer teto, a varredura para e o
resultado é marcado com \lstinline|truncated: true|. Cada ocorrência devolve
\lstinline|{ line, col, preview }|, e o caminho é relativo à raiz do projeto.

\section{Servidores de linguagem (LSP) para Verilog}
\label{sec:13-lsp}

A \tool{AURORA} integra dois servidores de linguagem distintos para
\textit{Verilog}/\textit{SystemVerilog}, com responsabilidades deliberadamente
divididas. Ambos falam JSON-RPC enquadrado por \lstinline|Content-Length| sobre
\textit{stdio}, por uma ponte de transporte escrita à mão — a \tool{AURORA} não
usa \textit{monaco-languageclient}; seu IPC já é próprio e um único \textit{backend}
de \textit{Verilog} precisa apenas de um transporte fino.

\subsection{Verible: o servidor sintático completo}

\file{main/lsp/verible\_lsp.js} gera um único processo de longa duração
\tool{verible-verilog-ls} (empacotado em
\path{components/Packages/verible/bin}), iniciado com
\lstinline|--lsp_enable_hover| e \lstinline|--rules_config_search| (para que o
\tool{Verible} procure, subindo a árvore, um \path{.rules.verible_lint} por
projeto). É o servidor \textbf{completo}: o \textit{renderer}
(\file{js/editor/lsp\_integration.js}) registra no \tool{Monaco} todo o conjunto
de provedores:

\begin{itemize}[leftmargin=*]
  \item \textbf{diagnósticos ao vivo} (\textit{lint} + sintaxe) como
        \textit{squiggles} e marcadores no painel \textit{Problems}, sob o dono
        \lstinline|verible|;
  \item \textbf{formatação} de documento (\textit{Format Document} /
        \lstinline|Shift+Alt+F|);
  \item \textbf{símbolos de \textit{outline}} (\textit{breadcrumbs} e visão
        \textit{Outline}), via \lstinline|provideDocumentSymbols|;
  \item \textbf{hover};
  \item \textbf{ir para definição} e \textbf{localizar referências}.
\end{itemize}

\begin{nota}
O \textbf{\textit{outline}} e o \textit{folding} simbólico vêm dos
\lstinline|documentSymbols| do \tool{Verible}, e \textbf{não} do
\textit{tree-sitter} (que é exclusivamente realce — ver
\autoref{sec:13-treesitter}). O \file{lsp\_integration.js} converte as
coordenadas LSP (base 0) para as do \tool{Monaco} (base 1) e mapeia
\lstinline|SymbolKind|/\lstinline|MarkerSeverity| entre os dois esquemas.
\end{nota}

O ciclo de vida é dirigido no nível do \textbf{modelo} do \tool{Monaco}: todo
modelo \path{.v}/\path{.sv} que aparece é \lstinline|didOpen|; suas edições são
debounced (350 ms) em \lstinline|didChange|; e o descarte dispara
\lstinline|didClose|. A sincronização é de documento inteiro (uma única mudança
sem \lstinline|range|). O \textit{bridge} é resiliente: se o servidor morre, ele
é respawnado na próxima requisição e os documentos abertos são re-enviados
(\lstinline|didOpen|) de forma transparente, de modo que os diagnósticos
continuam fluindo sem o \textit{renderer} perceber. Cada requisição
(\textit{format}/\textit{symbols}/\textit{hover}/...) tem \textit{timeout}
(\lstinline|REQUEST_TIMEOUT_MS| = 15 s) e, ao falhar, devolve um \textit{fallback}
vazio em vez de travar o editor.

\subsection{slang: o servidor semântico}

\file{main/lsp/slang\_lsp.js} gera o \tool{slang-server} (de
\textit{hudson-trading/slang-server}, em
\path{components/Packages/slang-server/bin}). Diferentemente do \tool{Verible}
(sintático e por arquivo), o \tool{slang} \textbf{elabora o projeto inteiro},
capturando erros semânticos que o \tool{Verible} não vê: identificadores não
declarados, incompatibilidades de tipo/porta, sinais não usados etc.

Pela divisão escolhida (\enquote{meio-termo}), o \textit{renderer}
(\file{js/editor/slang\_integration.js}) consome do \tool{slang} \textbf{apenas}:

\begin{itemize}[leftmargin=*]
  \item \textbf{diagnósticos semânticos} (sob o dono \lstinline|slang|,
        coexistindo com os marcadores \lstinline|verible|); e
  \item \textbf{completação} de símbolos (módulos, portas, membros de
        \textit{package}/\textit{struct}, macros), com caracteres de gatilho
        \lstinline| ` # . ( : [ |.
\end{itemize}

\textit{Hover}, definição, referências, \textit{outline} e formatação
permanecem com o \tool{Verible}. Como o \tool{slang} é acoplado ao
\textit{workspace} (indexa a árvore do projeto), o \textit{bridge} o inicia com
o diretório do projeto como \lstinline|rootUri| e o reinicia de forma
transparente quando o projeto muda (\lstinline|maybeRestartForProject|). Por
elaborar a cada mudança e poder ser ruidoso em projetos incompletos, o
\tool{slang} é \textbf{ligável/desligável}: \lstinline|window.AuroraSlang.toggle()|,
com estado persistido em \lstinline|localStorage| (\lstinline|slangEnabled|) e
sincronizado ao \textit{main}; há uma entrada na paleta de comandos para
alterná-lo. Ao desligar, os marcadores \lstinline|slang| são limpos e nenhuma
requisição é enviada.

\section{Realce sintático por tree-sitter}
\label{sec:13-treesitter}

O realce de sintaxe preciso é fornecido por \textit{tree-sitter} (via
\textbf{web-tree-sitter}, WASM), implementado em
\file{js/editor/treesitter\_highlight.js}. O \textit{parser} processa o
\textit{buffer}, a \textit{query} de \textit{highlights} da gramática produz
\textit{captures}, e o módulo os mapeia para \textbf{\textit{semantic tokens}}
do \tool{Monaco}, que sobrepõem o realce base (\textit{Monarch}) com cores
fiéis à gramática: nomes de módulo, instâncias, direção de portas, tipos,
macros etc.

\begin{nota}
O \textit{tree-sitter} na \tool{AURORA} faz \textbf{apenas realce} (semantic
tokens). Não provê \textit{folding} nem \textit{outline} — isso vem do LSP
(\autoref{sec:13-lsp}). As gramáticas cobertas são
\textit{Verilog}/\textit{SystemVerilog}, C e C++; \cmm{} e \textit{assembly}
não têm gramática e mantêm seus \textit{tokenizers} \textit{Monarch}.
\end{nota}

\subsection{Entrega das gramáticas WASM}

O \textit{renderer}, sob o \textit{sandbox} \lstinline|file://|, não consegue
buscar \path{.wasm} de forma confiável por URL. Por isso \file{main/treesitter/grammars.js}
serve os \textbf{bytes} crus de cada artefato pelo canal
\lstinline|treesitter:wasm|, e o \textit{renderer} os alimenta diretamente em
\lstinline|Parser.init({ wasmBinary })| e \lstinline|Language.load(bytes)|,
contornando todos os problemas de URL/\textit{fetch}/CSP. Apenas um conjunto fixo
de nomes lógicos é aceito (o \textit{renderer} não pode pedir caminhos
arbitrários):

\begin{table}[h]
\centering
\caption{Artefatos tree-sitter servidos por \file{grammars.js}.}
\label{tab:13-ts-artefatos}
\small
\begin{tabularx}{\linewidth}{Y Y}
\toprule
\textbf{Nome lógico} & \textbf{Arquivo} \\
\midrule
\lstinline|runtime|       & \path{web-tree-sitter.wasm} \\
\lstinline|systemverilog| & \path{tree-sitter-systemverilog.wasm} \\
\lstinline|c|             & \path{tree-sitter-c.wasm} \\
\lstinline|cpp|           & \path{tree-sitter-cpp.wasm} \\
\bottomrule
\end{tabularx}
\end{table}

Os artefatos ficam em \path{components/Packages/tree-sitter/} (baixados por
\file{download-tree-sitter-grammars.js} no \textit{bootstrap}). O canal
\lstinline|treesitter:status| informa se o \textit{runtime} e ao menos uma
gramática estão presentes.

\subsection{Detalhes de mapeamento}

As \textit{queries} (\path{js/editor/treesitter/queries/}) são incorporadas em
tempo de \textit{build} (\lstinline|?raw|); \textit{cpp} herda as \textit{queries}
de \textit{c}. Os nomes de \textit{capture} são mapeados para tipos de
\textit{token} semântico padrão do VS Code (\lstinline|CAPTURE_MAP|), de modo que
o tema ativo os colore sem configuração extra. Como as posições do
\textit{tree-sitter} são \textit{offsets} em \textbf{bytes} (UTF-8) e o
\tool{Monaco} usa colunas UTF-16, há conversão por linha (com caminho rápido
para ASCII), garantindo que comentários/\textit{strings} acentuados (por exemplo
em português) não desloquem o realce subsequente. Tudo é \textit{best-effort}:
se o \textit{runtime} ou os \textit{bytes} faltarem, o provedor devolve zero
\textit{tokens} e o realce \textit{Monarch} permanece, sem regressão. O
carregamento é \textbf{preguiçoso}: uma gramática só é buscada e compilada na
primeira vez que um \textit{buffer} daquela linguagem é realçado.

\section{Formatação com clang-format}
\label{sec:13-clang-format}

A formatação de C, C++ e \cmm{} é feita pelo \tool{clang-format} empacotado
(\path{components/Packages/clang-format/bin/clang-format.exe}), via
\file{main/format/clang\_format.js}. Diferentemente de um LSP, o
\tool{clang-format} é uma CLI de uma só passada: lê um \textit{buffer} no
\textit{stdin} e escreve o \textit{buffer} formatado no \textit{stdout}.

Por \lstinline|Shift+Alt+F|, o \textit{renderer}
(\file{js/editor/clang\_format\_integration.js}) registra um provedor de
formatação para \lstinline|c|, \lstinline|cpp| e \lstinline|cmm| e envia o
\textit{buffer} ao canal \lstinline|format:clang|. O \tool{Monaco} já despacha o
atalho para o provedor que casa com a linguagem do editor em foco — assim,
\textit{Verilog}/\textit{SystemVerilog} vão ao \tool{Verible} e
\lstinline|c|/\lstinline|cpp|/\lstinline|cmm| vão ao \tool{clang-format}. O
resultado é um \textbf{\textit{replace} de documento inteiro} (o
\tool{clang-format} devolve o \textit{buffer} completo, não um \textit{diff}).

O tratamento de linguagem usa \lstinline|-assume-filename|: o \cmm{} (um
subconjunto de C) é formatado com regras de C ao reivindicar um nome de arquivo
\path{.c}; C/C++ usam o caminho real, de modo que o \tool{clang-format} detecte o
dialeto \textbf{e} encontre, subindo a árvore, um \path{.clang-format} de
projeto (executado com \lstinline|-style=file -fallback-style=LLVM|). Há
\textit{timeout} de 15 s contra travamentos, verificação de \textit{allowlist}
de binário (\lstinline|isAllowed|, a mesma do executor do \textit{toolchain}) e
comportamento \textit{best-effort}: se o binário faltar ou houver erro,
\lstinline|format()| resolve \lstinline|null| e o \textit{buffer} fica
intocado.

\section{Paleta de comandos}
\label{sec:13-paleta}

A paleta de comandos (\file{js/ui/command\_palette.js}) é uma superfície única,
orientada a teclado, para as ações espalhadas pela barra de ferramentas e
menus. Abre com \lstinline|Ctrl/Cmd+Shift+K| (principal) ou
\lstinline|Ctrl/Cmd+Shift+P|; \lstinline|Esc| fecha, setas navegam,
\lstinline|Enter| executa. O \lstinline|Ctrl+K| puro é reservado para o painel
de IA e por isso não é vinculado. A captura é feita na fase de captura com
\lstinline|stopPropagation|, de forma que a paleta abra mesmo sobre um editor
\tool{Monaco} em foco (que liga \lstinline|Ctrl+Shift+K| a \enquote{apagar
linha}) sem disparar aquele comando.

A \textbf{visão} é o componente Lit \lstinline|<aurora-command-palette>| (Shadow
DOM + \textit{tokens} semânticos); o módulo é o dono do \textbf{registro}
(dados puros), da \textbf{pontuação} \textit{fuzzy} e da lógica de execução. Os
comandos preferem a API pública (\lstinline|window.AuroraAPI| / controlador da
árvore de arquivos) e, na ausência dela, clicam o botão de barra existente por
\lstinline|id| — assim um comando faz exatamente o que o botão faz (inclusive
ser \textit{no-op} quando o botão está desabilitado), sem lógica duplicada.

O registro \lstinline|COMMANDS| é agrupado em \lstinline|Compile|,
\lstinline|Project|, \lstinline|View|, \lstinline|Tools| e \lstinline|Dev| (a
ordem de exibição segue \lstinline|GROUP_ORDER|). Inclui, entre outros:
compilar \cmm, sintetizar \textit{Verilog}, gerar \textit{waveform}, execução
rápida (\tool{Verilator}), abrir o \tool{PRISM}, novo/abrir projeto,
configurações da \tool{AURORA}, alternar o \tool{slang} e alternar o
\textit{overlay} de \textit{jank} (este último por \textit{import} dinâmico, sem
custo de \textit{boot}). A função \lstinline|scoreCommand| pontua por
subsequência: cada termo da consulta precisa aparecer no \textit{haystack}
(título + \textit{keywords} + grupo); acertos no título (e prefixos) valem mais
que acertos em \textit{keywords}.

\section{Internacionalização (i18n)}
\label{sec:13-i18n}

A camada de internacionalização vive em \file{js/i18n/i18n.js}, com dois
catálogos em \path{locales/}: \path{en.json} e \path{pt.json}. São dois
\textit{locales} suportados — inglês e português do Brasil — e o padrão é
\lstinline|pt|. O DOM declara associações por dois atributos principais:

\begin{itemize}[leftmargin=*]
  \item \lstinline|data-i18n="caminho.da.chave"| $\rightarrow$
        \lstinline|element.textContent|;
  \item \lstinline|data-i18n-tooltip| $\rightarrow$ \lstinline|data-tooltip|;
        \lstinline|data-i18n-placeholder| $\rightarrow$ \lstinline|placeholder|;
        \lstinline|data-i18n-aria-label| $\rightarrow$ \lstinline|aria-label|;
        \lstinline|data-i18n-title| $\rightarrow$ \lstinline|title|.
\end{itemize}

A função \lstinline|applyDOM()| percorre o documento no \textit{boot} e a cada
troca de \textit{locale}, resolvendo cada associação por \lstinline|t(key)|. A
mesma \lstinline|t()| é exposta como \lstinline|window.t| para chamadas em JS
(mensagens do terminal, UI dinâmica). Há ainda um preenchimento automático de
\textit{tooltip} por \lstinline|id|: qualquer elemento com \lstinline|id="X"|
recebe \lstinline|data-tooltip| de \lstinline|tooltip.extended.X| quando a chave
existe.

\subsection{Resolução, \textit{fallback} e interpolação}

O \textit{locale} ativo é guardado em \lstinline|localStorage| sob
\lstinline|aurora-locale| (a chave legada \lstinline|aurora-yanc-lang| é migrada
silenciosamente na primeira leitura). A resolução tenta o \textit{locale} ativo;
se a chave faltar, recai em \lstinline|en|; se ainda faltar, devolve o próprio
caminho da chave — uma falha \enquote{barulhenta} que evidencia strings sem
tradução durante o desenvolvimento. A interpolação resolve marcadores
\lstinline|{{nome}}| no \textit{template}; parâmetros ausentes viram string
vazia (sem exceção).

Cada troca de \textit{locale} dispara o evento
\lstinline|aurora:locale-changed|, ao qual componentes que renderizam strings
imperativamente (terminal, notificações) reagem para re-renderizar. Por
compatibilidade com o compilador \tool{YANC}, \lstinline|window.getYancLang| é
um \textit{alias} de \lstinline|getLocale()|, de modo que a injeção do
\textit{flag} \lstinline|-pt|/\lstinline|-en| do compilador segue inalterada — a
preferência de idioma é, intencionalmente, unificada entre UI e compilador. O
controle visual é o \lstinline|<select>| de idioma em
\textbf{Configurações $\rightarrow$ Idioma}, ligado por
\file{js/ui/locale\_toggle.js}.

\section{Painel de configurações}
\label{sec:13-config}

O painel de configurações é montado por \file{js/processors/aurora\_settings.js},
acionado pelo botão \lstinline|aurora-settings| (e pelo comando homônimo na
paleta). Ele abre um modal (\lstinline|settings-modal|) com os controles de
preferências. O estado é persistido em \lstinline|localStorage|:

\begin{itemize}[leftmargin=*]
  \item \lstinline|aurora-settings| — as preferências gerais (por exemplo,
        \textit{tooltips} habilitados);
  \item \lstinline|aurora-shortcuts| — os atalhos de teclado personalizados;
  \item \lstinline|aurora-ai-trust-external-links| — chave compartilhada com a
        caixa de \enquote{confiar em links externos} do chat de IA, de modo que
        as duas fiquem ligadas.
\end{itemize}

Os atalhos padrão (\lstinline|defaultShortcuts|) cobrem ações como novo arquivo
(\lstinline|Ctrl+N|), compilar tudo (\lstinline|Ctrl+Shift+B|), fechar/reabrir
aba, salvar/salvar tudo e abrir configurações. Os rótulos dos atalhos são
resolvidos em tempo de renderização via chaves de i18n
(\lstinline|SHORTCUT_LABEL_KEYS|), não armazenados — assim uma troca de
\textit{locale} atualiza a UI sem tocar nos atalhos persistidos. O painel também
expõe o \textit{select} de idioma (\autoref{sec:13-i18n}) e o controle de
\textit{tooltips} (aplicado imediatamente por \lstinline|setTooltipsEnabled|).

\section{Atualizador automático (do ponto de vista do usuário)}
\label{sec:13-updater}

O \textit{auto-updater} vive em \file{main/updater.js} e é construído sobre
\textbf{electron-updater}, consultando os \textit{releases} do repositório
\repo{sapho} no GitHub. Os detalhes de \textit{release} e empacotamento são
tratados em \autoref{cap:16-release}; aqui descreve-se a experiência do usuário.

O fluxo (somente em produção — pulado em modo de desenvolvimento) é:

\begin{enumerate}[leftmargin=*]
  \item Cerca de 6 segundos após a janela principal aparecer, a \tool{AURORA}
        checa silenciosamente por atualizações
        (\lstinline|autoUpdater.autoDownload = false|).
  \item Em \lstinline|update-available|, abre-se uma janela de atualização
        personalizada (\path{html/update-notification.html}) mostrando o
        \textit{changelog} bilíngue e a opção de baixar — nunca um diálogo
        nativo. Se o \textit{electron-updater} não trouxer as notas, elas são
        buscadas diretamente na API do GitHub (\lstinline|fetchReleaseNotes|).
  \item Ao iniciar o \textit{download}, o evento \lstinline|download-progress|
        alimenta uma barra de progresso real (percentual, MB transferidos,
        velocidade, ETA) na mesma janela.
  \item Em \lstinline|update-downloaded|, a janela muda para o estado
        \enquote{Reiniciar e instalar}.
  \item \lstinline|quitAndInstall(false, true)| instala e relança diretamente
        na nova versão.
\end{enumerate}

Uma verificação manual (\enquote{checar atualizações}, via menu/configurações/
\textit{About}) usa \lstinline|checkForUpdates(true)| e, se nada for encontrado,
exibe um diálogo \enquote{você está atualizado}; a checagem silenciosa de
inicialização suprime esse aviso. Erros são traduzidos para mensagens curtas e
legíveis por \lstinline|friendlyError| (falha de rede, verificação de
assinatura, falta de espaço em disco, permissão negada). Há ainda a
\textbf{detecção pós-atualização}: a versão em execução é gravada em
\path{userData/aurora-version.json} a cada \textit{boot}; no \textit{boot}
seguinte, uma divergência entre a versão armazenada e a atual indica que o
usuário acabou de atualizar — o \textit{renderer} lê esse estado uma vez
(\lstinline|updates:post-update-status|) e mostra um \textit{toast} de
confirmação. Instalações novas não disparam falso \enquote{você foi
atualizado}, e o sinalizador é de uso único (uma segunda janela na mesma sessão
não repete o \textit{toast}).

\section{Overlay de desenvolvimento (Jank Overlay)}
\label{sec:13-jank}

Para diagnóstico de desempenho durante o desenvolvimento, há o \textit{Jank
Overlay} em \file{js/dev/jank\_overlay.js}, um HUD que mede e exibe em tempo
real:

\begin{itemize}[leftmargin=*]
  \item \textbf{FPS} instantâneo (último quadro);
  \item \textbf{p99} do tempo de quadro sobre uma janela deslizante de 300
        quadros;
  \item \textbf{taxa de \textit{jank}}: percentual de quadros cujo tempo excede
        o dobro do orçamento de 165 Hz --- isto é, $2 \times 6{,}06 \approx
        12$ ms (a constante \lstinline|BUDGET_MS| vale $1000/165 \approx 6{,}06$ ms);
  \item contagem de \textbf{\textit{longtasks}} (via \lstinline|PerformanceObserver|,
        tarefas $>$ 50 ms);
  \item aproximação de \textbf{TTI} (\textit{time to interactive}), a partir da
        marca \lstinline|aurora-interactive| ou do \lstinline|domInteractive| da
        entrada de navegação.
\end{itemize}

O \textit{overlay} tem \textbf{custo zero quando inativo}: o laço de
\lstinline|requestAnimationFrame| e o \lstinline|PerformanceObserver| só rodam
enquanto ele está visível. É alternado pela paleta de comandos
(\enquote{\textit{Toggle Jank Overlay}}, grupo \lstinline|Dev|), e o
\textit{import} dinâmico do módulo evita qualquer custo de inicialização da IDE.
Os valores são coloridos por limiares (verde/âmbar/vermelho) para leitura
rápida.

\section{Síntese}
\label{sec:13-sintese}

As ferramentas de produtividade da \tool{AURORA} compõem uma camada coesa em
torno do editor, conectando processos privilegiados do \textit{main} a
integrações no \tool{Monaco} e à interface. O terminal de saída estrutura e
torna navegável o que as ferramentas do \textit{toolchain} emitem; o controle de
versão embutido (sobre \textbf{simple-git} e \textbf{diff2html}) cobre o ciclo
completo de \lstinline|git| com autenticação GitHub cifrada; a busca global é um
varredor próprio sem dependências; o par \tool{Verible} (sintático completo) +
\tool{slang} (semântico, opcional) entrega diagnósticos, formatação,
\textit{outline}, \textit{hover}, navegação e completação; o
\textit{tree-sitter} acrescenta realce fiel à gramática; o \tool{clang-format}
formata C/C++/\cmm; e paleta de comandos, i18n, configurações, atualizador e
\textit{overlay} de \textit{jank} fecham a experiência. O fio condutor é o
\textit{best-effort}: cada recurso degrada graciosamente para um \textit{no-op}
quando seu binário ou dependência não está presente, sem nunca comprometer a
edição.
